El presente artículo científico presenta en detalle el diseño, la implementación y la 
validación de un sistema de gestión de tickets, ofreciendo una solución tecnológica para 
optimizar el flujo de trabajo en instituciones donde la atención a incidencias y 
requerimientos se maneja de forma manual o poco estructurada. El proyecto propone una 
arquitectura híbrida que integra un portal web y una aplicación móvil Android, ambos 
conectados a una base de datos centralizada, con módulos de seguridad, notificaciones en 
tiempo real y funcionalidades para la clasificación automática de tickets. El documento 
describe las etapas de desarrollo y se fundamenta en ocho artículos seleccionados que 
abordan frameworks de desarrollo modernos, metodologías ágiles y patrones 
arquitectónicos relevantes. Estos estudios incluyen aportes sobre frameworks PHP como 
Laravel y Symfony [1], [7], bibliotecas front-end como React y Vue [2], [4], [5], React 
Native para aplicaciones móviles [3], y enfoques metodológicos basados en Scrum y en 
el conjunto de metodologías ágiles [6], [9]. A partir de este contraste, se evidencian las 
ventajas de integrar teoría y práctica en un mismo proyecto, validando que la 
incorporación de frameworks modernos y enfoques ágiles permite obtener aplicaciones 
escalables, seguras y con una alta orientación hacia la experiencia de usuario.